\begin{theorem}[Regular Cauchy minimum-modulus sibling route]
\label{thm:regular-cauchy-min-modulus-sibling-route}
Every $\RegularCauchyMinModulusUp$ read factors through the $\CauchyModulusUp$ carrier of \autoref{ch:concrete-instances-cauchymodulus-namecert} and the $\RegularCauchyModulusNormalizerUp$ normalizer of \autoref{ch:concrete-instances-regular-cauchy-modulus-normalizer-namecert} before exposing the minimum modulus row. The route first displays the source modulus request, then normalizes the regular-Cauchy threshold schedule, and only then reads the finite minimum row together with its $\BHist$, $\hsame$, $\Cont$, $\Pkg$, and local $\NameCert$ support.
\end{theorem}
\begin{proof}
Start with the carrier of \autoref{def:regular-cauchy-min-modulus-carrier}. The accepted packet displays source schedules, the meet row, the minimum row, stream windows, regular-rational readback, dyadic tolerance data, the Real seal, transport, replay, provenance, and local naming.

Read the $\CauchyModulusUp$ sibling first. It supplies the carrier-level modulus request whose finite thresholds are visible before the minimum row can be inspected.

Next read the $\RegularCauchyModulusNormalizerUp$ sibling. Its normalizer keeps the regular-Cauchy threshold schedule on the same displayed carrier route, so the minimum row is not a detached host-level comparison.

The exposed minimum modulus row is therefore downstream of both sibling packets. Since the route remains inside the displayed $\BHist$, $\hsame$, $\Cont$, $\Pkg$, and local $\NameCert$ rows, it exports no quotient completion, selected limit, countable-choice selector, or ambient Real completeness principle.
\end{proof}
